% ============================================
% DOC-13
% Gerado de molambudos.md — build_miolo.py
% ============================================
\fragtitle{doc}{D}{DOC}{13}{Diário de Oliveira III}
\markboth{DOC-13 --- Diário de Oliveira III}{}

\noindent \textit{A última sessão com Joaquim}

\noindent \textbf{Dr. Heitor Oliveira --- Caderno Pessoal}
\textbf{12 de junho de 1979}

\vspace{0.5em}\noindent\rule{\textwidth}{0.5pt}\vspace{0.5em}

\noindent \textit{Nota do arquivista: Esta é a última entrada escrita pelo Dr. Oliveira no dia da morte do paciente 1.260. A caligrafia está visivelmente alterada. Algumas palavras estão riscadas. Outras foram escritas com tanta força que furaram o papel.}

\vspace{0.5em}\noindent\rule{\textwidth}{0.5pt}\vspace{0.5em}

\noindent ELE MORREU.

\noindent Estou escrevendo no corredor do hospital, com as mãos sujas de sangue e terra. Não, não é sangue. É tinta. A tinta do diário vazou. Ou eu sujei a mão na parede. Não sei.

\noindent O paciente 1.260 morreu hoje.

\noindent Estava na sessão. Ele estava calmo. Mais calmo do que nunca. Olhou para mim com o olho amarelo --- o mesmo olho que agora está no meu rosto --- e disse:

\noindent \textit{\textquotedblleft{}O senhor está diferente.\textquotedblright{}}

\noindent Perguntei como.

\noindent \textit{\textquotedblleft{}O olho. Está clareando. Já começou.\textquotedblright{}}

\noindent Toquei meu rosto. Não senti diferença. Mas ele viu.

\noindent Ele viu porque ele já passou por isso.

\noindent Ele viu porque ele é o espelho.

\noindent \textit{\textquotedblleft{}O que acontece agora?\textquotedblright{}}, perguntei.

\noindent \textit{\textquotedblleft{}Acontece que eu vou morrer.\textquotedblright{}}

\noindent \textit{\textquotedblleft{}Morrer? Por quê?\textquotedblright{}}

\noindent \textit{\textquotedblleft{}Porque a criatura precisa de espaço. Não cabe mais dentro de mim. Eu sou velho. Fraco. Ela precisa de um corpo novo. O senhor.\textquotedblright{}}

\noindent \textit{\textquotedblleft{}Não.\textquotedblright{}}

\noindent \textit{\textquotedblleft{}Já é, doutor. O senhor leu o diário. O senhor sonhou com a vala. O senhor sentiu o cheiro. O senhor está amarelado. Já é.\textquotedblright{}}

\noindent Eu neguei. Disse que ia internar ele. Que ia queimar o diário. Que ia impedir.

\noindent Ele sorriu.

\noindent \textit{\textquotedblleft{}Pode queimar. A história não está no papel. Está na sua cabeça. A criatura não é feita de tinta. É feita de palavras. E as palavras estão dentro do senhor agora.\textquotedblright{}}

\noindent Ele estendeu a mão.

\noindent \textit{\textquotedblleft{}Obrigado.\textquotedblright{}}

\noindent \textit{\textquotedblleft{}Obrigado pelo quê?\textquotedblright{}}

\noindent \textit{\textquotedblleft{}Por me substituir. Estou cansado. 72 anos carregando ela. Agora posso descansar.\textquotedblright{}}

\noindent Ele fechou os olhos.

\noindent E morreu.

\noindent Na minha frente.

\noindent Na cadeira da sessão.

\noindent O olho amarelo apagou.

\noindent Mas, no mesmo instante, o MEU olho direito acendeu.

\noindent Senti uma dor. Uma queimadura. Corri para o espelho do banheiro.

\noindent O olho direito está dourado.

\noindent Brilhando.

\noindent Igual ao dele.

\noindent OLHO IGUAL AO DELE.

\vspace{0.5em}\noindent\rule{\textwidth}{0.5pt}\vspace{0.5em}

\noindent O diário está na minha mão.

\noindent Ele está quente.

\noindent As páginas estão se movendo sozinhas.

\noindent Estão se enchendo de letras. A minha letra. Mas eu não estou escrevendo.

\noindent A mão está escrevendo sozinha.

\noindent Está escrevendo:

\noindent \textit{\textquotedblleft{}O CICLO RECOMEÇA. PACIENTE 1.261. HEITOR OLIVEIRA.\textquotedblright{}}

\noindent Não.
NÃO.

\noindent NÃO É VERDADE.

\noindent ISTO É UM DELÍRIO. ISTO É SUGESTÃO. ISTO É---

\noindent \textit{tô escrevendo com a mão esquerda agora. a direita não para. ela escreve sozinha. o diário tá cheio de palavras que eu não escrevi. mas são minhas. a letra é minha. mas não fui eu.}

\noindent \textit{ele tá escrevendo através de mim.}

\noindent \textit{a criatura.}

\noindent \textit{ela tá escrevendo.}

\noindent \textit{a mão.}

\noindent \textit{não para.}

\noindent \textit{NÃO PARA.}

\noindent \textit{NÃ---}

\vspace{0.5em}
\noindent\textit{\textbf{↪ Rotas:} \rota{DOC-14} \rotasep \rota{MEM-16} \rotasep \rota{LUC-04}}